\chapter{Descrizione del tirocinio}
\label{chap:descrizione-tirocinio}
\noindent
\textit{Nel seguente capitolo spiego il ruolo dei tirocini all'interno di Bluewind, contestualizzandoli in relazione alla loro visione di innovazione.
Vado poi a descrivere l'attività di tirocinio che mi è stata proposta, quali motivazioni hanno spinto Bluewind a ideare quest'iniziativa e con quali obiettivi.
Infine, do le mie, personali, motivazioni che mi hanno portato a scegliere questa proposta di stage.}

\section{Gli stage all'interno del piano aziendale}
\iffalse
In questa sezione parlo di come gli stage si integrino all'interno del piano aziendale e quali obiettivi l'azienda cerca di perseguire con essi.
\fi
\noindent
Come già anticipato nella sezione \ref{sect:prop-innov}, la continua necessità e ambizione di \textit{Bluewind} di rimanere al passo con il progresso tecnologico, porta l'azienda a riservare tempo e personale per inseguire tale obiettivo.
Non è però scontato riuscire ad avere sempre a disposizione i mezzi per fare ciò, soprattutto visto l'avanzamento sempre più rapido del settore informatico.\\

\noindent
Per questo motivo l'azienda cerca di sfruttare appieno l'opportunità rappresentata dai tirocini, lavorando a stretto contatto con le università.
I tirocini universitari permettono infatti a \textit{Bluewind} di esplorare con più facilità le tematiche di suo interesse.
Vengono raccolti nel tempo spunti e idee su tecnologie, strumenti e progetti che potrebbero essere di impatto per la crescita dell'azienda, e partendo da essi si vanno a creare delle proposte di tirocinio con lo scopo valutarne l'adozione.
Le proposte di tirocinio possono dunque avere sia come obbiettivo la ricerca dei suddetti temi, sia la creazione di un \gls{poc}, un prodotto dimostrativo che serve a valutare la fattibilità e l'utilità di un progetto prima che ne si cominci la realizzazione vera e propria.
A partire dai risultati ottenuti, l'azienda può poi decidere se è opportuno proseguire a esplorare o meno la tecnologia in questione, avendo anche, nel caso del \textit{PoC}, una base concreta dalla quale poter partire per lo sviluppo di un progetto.\\

\noindent
Gli \textit{stage} rappresentano inoltre un'opportunità per l'azienda nel merito delle assunzioni.
I tirocini permettono infatti all'azienda di seguire da vicino ragazzi giovani e prossimi alla laurea, valutando con concretezza la possibilità di inserirli all'interno della propria attività.
Questa modalità di \textit{scouting}, che ha ,negli anni, portato all'assunzione di diverse persone, ha oltretutto il vantaggio di agevolare l'ambientamento da parte dei nuovi dipendenti, in quanto essi hanno già avuto modo di conoscere quella realtà.

\section{La spinta normativa}
\noindent
Molte scelte e attività all'interno di \textit{Bluewind} sono fortemente influenzate da fatto che nel loro settore vi sono diverse normative stringenti a cui le aziende devono sottostare.
In questa sezione vorrei quindi spendere due parole per rendere nota la necessità dell'azienda di aderire a certi \textit{standard}.
Ritengo che sia un punto abbastanza importante quantomeno da introdurre per dare a chi legge una visione più chiara sul alcune dinamiche.

\section{Il progetto e la sua finalità}
\noindent
In questa sezione parlo del progetto propostomi, da cosa nasce questa necessità e quale ruolo può avere all'interno del contesto aziendale.

\section{Gli obiettivi del progetto}
\noindent
In questa sezione descrivo gli obiettivi prefissati con l'azienda da raggiungere durante l'attività di tirocinio.

\section{Le ambizioni personali}
\noindent
In questa sezione parlo delle ragioni che mi hanno spinto a considerare questa proposta di stage, in relazione anche ai miei piani futuri di carriera.
Oltre a ciò specificherò quali sono gli obiettivi personali che mi sono posto di raggiungere con quest'esperienza in maniera quanto più quantitativa possibile.
